\subsection{A nonzero theorem at the new bootstrap}
The new complete source frame has register vector
\[
 (\mathsf{np},\mathsf{tw},\mathsf{ws},\mathsf{ac},d,h,s_1,s_2,\mathsf{pl})
 =(4,3,0,0,3,0,0,0,0).
\]
Here native code $3=E(0,0)$ is an equality theorem; native $I(0,0)=0$.
In particular, neither the counter cutoff four nor the nonzero additive
parameter invalidates the enumeration argument.

\begin{lemma}[Admissibility of the new logical seed]
The unchanged native enumerator, started from the displayed clean root
vector, halts if and only if the retained calculus $C_{222}$ derives its
falsum code $F=1$.
\end{lemma}
\begin{proof}
First, $C_{222}\vdash E(x,x)$ for every natural $x$.
Put $r=R(x,x)$ and $X=E(r,x)$. The merged equality axiom is
$X\to(X\to Q_rE(x,x))$; contraction gives
$X\to Q_rE(x,x)$. Existential monotonicity, derived by contraposition,
closed generalization and $\rho_2$, gives
$\exists_rX\to\exists_rQ_rE(x,x)$.
Apply $\rho_1(r,x)$ and then $\rho_6(r,E(x,x))$.
Specializing $x=0$ proves the new current formula $3$.

Thus the usual invariant holds initially: the current formula and cache
are theorems and every decoded stack entry is a theorem. The empty stack
is the all-zero stream, and zero is truth. The invariant imposes no
condition on the candidate counter or additive parameter. Every command
preserves it, and a halting NIP tests a cached falsum. This proves
soundness without assuming consistency.

For completeness, compile any derivation of $F$ to a finite block $p$
whose contract is universal in the current formula $T$ and underlying
stack $S$: from $(T,S,0,0)$ it either halts or returns $(F,S,0,0)$.
For any natural $U$, the command
$\operatorname{START}_U=\operatorname{AX}(0,1,1,U)$ resets the additive
parameter and cache to zero. Hence
\[
 \operatorname{START}_U\mathbin{+\!+}p\mathbin{+\!+}
 [\operatorname{DROP}(0),\operatorname{NIP}(0),\operatorname{NIP}(0)]
\]
halts from every working tuple. If $R$ is the encoded pointer of the fixed
nonempty suffix, its complete candidate number is
\[
 N_U=P\bigl(3,P(1,P(1,P(U,R)))\bigr)\geq U.
\]
Choose $U=4$. Nonzero proof-list pointers strictly decrease on every
nonhalting command; when the pointer is zero, the candidate counter
increments. Consequently the run eventually executes candidate $N_U$
unless it already halted. This establishes completeness above cutoff four.
All constructions are finite syntactic transformations and natural-number
inductions formalizable in ZFC.
\end{proof}

The existing $\texttt{L2Root219.toAS}$ is precisely
$[\mathsf{np},\mathsf{tw},\mathsf{ws},0,d,h,0,0,\mathsf{pl}]$;
its root-step correspondence is universal in the enumerator state.
The accepted K218 source equivalence is universal in the register vector.
Thus this lemma fits the unchanged source directly. A complete binary
bootstrap and the primitive-runtime bridge remain separate physical
obligations; the lemma does not substitute for either of them.
